% The abstract of Paper 1. Learned Publishing Opinion pieces have none, so the journal
% build leaves it out; the preprint build (make paper1-arxiv, \preprintmode) pulls it in
% after \maketitle. It does not count against the journal word budget.

\begin{abstract}
Generative AI tools now help write a growing share of academic papers.
Checking their citations, that each cited work exists and says what it
is cited for, has received attention.  A paper can pass that check and
still fail: a claim can rest on real evidence that does not carry it,
or be worded more certainly than its support allows.  Reporting
guidelines, such as those collected by the EQUATOR (Enhancing the
QUAlity and Transparency Of health Research) Network, help authors
report a study completely; whether the reasoning from evidence to claim
holds is a separate question, most exposed in papers that have only
reasoning to offer, and a documented search found no consensus
checklist that checks such reasoning step by step.  This paper proposes
a claim register, built on Toulmin's layout of arguments and Whetten's
account of a theoretical contribution: each claim is recorded with what
it rests on, each kind of grounds gets its own check, the step from
grounds to claim is checked for relevance and sufficiency, and the
wording of each claim is tied to how well it is supported.  Each part
is already in use, in safety engineering, scholarly data publishing,
educational measurement and the calibrated reporting of certainty; none
of the approaches compared here combines them for academic papers.
Applied to this paper, the register leaves three of its six central
claims at the Emerging tier, for named reasons: the search had one
screener, and the records were kept by the author's own tooling.
Publishing a register can make a paper's reasoning visible and
traceable; whether it helps authors and reviewers find weak steps has
not been tested.  Authors and journals could publish registers as
optional supplementary material, and trials should test them before
wider adoption.

\medskip\noindent
\textbf{Keywords:} academic writing, generative AI, verification,
argumentation, reporting guidelines, non-empirical research
\end{abstract}
